% Generated by build/sync_qft_soln.py from committed sources only.
% Source repository: https://github.com/Luca-Yucheng-Jin/QFT-soln
% Source commit: 07d7fe95fd5cf5b26d38a16ee535a04a925277b2
\section{Renormalization (PSI QFT II, Tutorial 6)}

\noindent\textit{Question prompts paraphrased from the official
\href{https://psi-online.perimeterinstitute.ca/courses/take/quantum-field-theory-ii-student/pdfs/20776391-module-5-tutorial}{Perimeter PSI Online sheet};
the equations and notation follow the source.}

\subsection{RG Flow in Three Dimensions}

Consider massless Euclidean $\phi^4$ theory in three spacetime dimensions,
\begin{equation}
    S_E[\phi]
    =\int d^3x\,
    \left[
        \frac{1}{2}(\partial\phi)^2
        +\frac{g}{4!}\phi^4
    \right].
\end{equation}

\begin{enumerate}[label=\alph*)]
    \item Determine the mass dimension of the coupling $g$.
    \begin{solution}
        $g$ has mass dimension $4-d = 1$.
    \end{solution}

    For two-to-two scattering with incoming momenta $p_1,p_2$ and outgoing
    momenta $p_3,p_4$, define the Mandelstam variables
    \begin{equation}
        s=(p_1+p_2)^2,
        \qquad
        t=(p_1-p_3)^2,
        \qquad
        u=(p_1-p_4)^2.
    \end{equation}
    In the conventions of the sheet,
    \begin{equation}
        s+t+u=-4m^2=0,
    \end{equation}
    where the final equality uses $m=0$.

    \item Write the four-point irreducible Green function
    $\overline\Gamma_{(4)}$ through one loop in terms of
    $p_1,p_2,p_3,p_4$. Set up the loop integrals, but do not evaluate them yet.
    \begin{solution}
        \begin{equation}
            \begin{aligned}
                -\overline\Gamma_{(4)}
                &= -g + \frac{(-g)^2}{2}\int d^dk \frac{1}{k^2(k+p_1+p_2)^2} + (p_2 \leftrightarrow -p_3) + (p_2 \leftrightarrow -p_4) + \cdots.
            \end{aligned}
        \end{equation}
    \end{solution}

    \item In four dimensions the corresponding integrals are ultraviolet
    divergent. Show, by power counting or direct analysis, that the
    three-dimensional loop integrals are ultraviolet finite. No UV regulator
    is needed here.
    \begin{solution}
        For $d=4$, the integral is logarithmic divergent. For $d=3$, it remains finite.
    \end{solution}

    \item Show that
    \begin{equation}
        \overline\Gamma_{(4)}(s,t,u)
        =g-\frac{g^2C}{2}
        \left[
            \sqrt{s^{-1}}+\sqrt{t^{-1}}+\sqrt{u^{-1}}
        \right],
        \qquad C>0.
    \end{equation}
    Determining the precise value of $C$ using Feynman parameters is optional.
    \begin{solution}
        \begin{equation}
            \begin{aligned}
                \overline\Gamma_{(4)} &= g - \frac{g^2}{2}\left( 2\pi\int ^1_{-1}d\cos\theta \int^\infty_0 dk \frac{1}{k^2 + s +2\sqrt{s}\cos\theta} +(s\leftrightarrow t) + (s\leftrightarrow u)\right)\\
                &=g - \frac{g^2}{2}\left( \frac{\pi}{\sqrt{s}} \int^\infty_0\frac{dk}{k} \log\left(\frac{k^2 + s+2\sqrt{s}}{k^2 + s-2\sqrt{s}}\right) +(s\leftrightarrow t) + (s\leftrightarrow u)\right)\\
                &=g - \frac{g^2}{2}\left( \frac{2\pi}{\sqrt{s}} \int^\infty_0\frac{dk}{k} \log\left(\frac{k + \sqrt{s}}{k-\sqrt{s}}\right) +(s\leftrightarrow t) + (s\leftrightarrow u)\right)\\
                &=g - \frac{g^2}{2}\left( \frac{2\pi}{\sqrt{s}} \int^\infty_0\frac{dk}{k} \log\left(\frac{k + 1}{k-1}\right) +(s\leftrightarrow t) + (s\leftrightarrow u)\right)\\
                &=g - \frac{g^2}{2}2\pi\int^\infty_0\frac{dk}{k} \log\left(\frac{k + 1}{k-1}\right)(\sqrt{s^{-1}}+\sqrt{t^{-1}}+\sqrt{u^{-1}}),
            \end{aligned}
        \end{equation}
        where we have rescaled $k \to k/\sqrt{s}$ in the second last equality.
    \end{solution}

    Although the symmetric choice $s=t=u=\mu^2$ is incompatible with the
    preceding Mandelstam relation, adopt it as the simplified subtraction point
    used in the sheet and define
    \begin{equation}
        h_R
        =\mu^{-1}\overline\Gamma_{(4)}
        (s=t=u=\mu^2).
    \end{equation}
    Hence,
    \begin{equation}
        h_R
        =\mu^{-1}
        \left[
            g-\frac{3}{2}g^2\mu^{-1}C
        \right].
    \end{equation}

    \item At an arbitrary energy $E$, define
    \begin{equation}
        h_{\mathrm{eff}}(E)
        =E^{-1}\overline\Gamma_{(4)}(s=t=u=E^2),
    \end{equation}
    together with
    \begin{equation}
        \beta(h_{\mathrm{eff}})
        =E\frac{dh_{\mathrm{eff}}}{dE}.
    \end{equation}
    Show, through order $g^2$, that
    \begin{equation}
        \beta(h)
        =-h+\frac{3}{2}Ch^2+\text{higher-order terms}.
    \end{equation}
    \begin{solution}
        \begin{equation}
            \beta(h) = E\left(-\frac{1}{E^2}\left(g - \frac32g^2E^{-1}C\right) + \frac{1}{E}\cdot \frac32 g^2\frac{1}{E^2}C\right) = -h + \frac32 Ch^2 + \mathcal{O}(g^3).
        \end{equation}
    \end{solution}

    \item Plot $\beta(h)$ as a function of $h$.

    \item Suppose $h$ is small and positive at some finite energy. Show that
    $h\to0$ as $E\to\infty$, and identify this zero as the ultraviolet fixed
    point of the RG flow.
    \begin{solution}
        Solving the RG equation, we find
        \begin{equation}
            h(E) = \frac{1}{\frac{3}{2}C-\frac{h_0}{\frac32Ch_0-1}\frac{E}{E_0}}.
        \end{equation}
        Clearly, as $E\to \infty$, $h \to 0$. Note that $\beta(0) = 0$, and hence this is a UV fixed point.
    \end{solution}

    \item Starting again from small positive $h$, show that the flow approaches
    a finite value $h_0$ as $E\to0$, where $h_0$ is the nonzero root of
    $\beta(h)$. Interpret $h_0$ as an infrared fixed point.
    \begin{solution}
        As $E \to 0$, $h \to \frac{2}{3C}$. Note that $\beta(\frac{2}{3C}) =0$, and hence this is an IR fixed point. This suggests the existence of the Wilson-Fisher fixed point. Note that at this point, the theory become strongly coupled, since all higher order terms become relevant.
    \end{solution}
\end{enumerate}
